\begin{cnabstract}[自动调制识别；RadioML2016.10A；时频特征；多视图融合；CLDNN；信号域增强；配对统计检验]
低信噪比下的自动调制识别中，多视图时频融合是否真正有效尚需实证检验。本文以 \datasetname{} 固定划分 \splitname{} 与三个训练种子建立 9 模型 \(\times\) 3 种子 \(=27\) 运行单元的协议化主实验，并以独立证据标签登记四类扩展实验与一个提案模型。主实验中 CLDNN 取得最高整体准确率 0.6129、宏平均 F1 0.6332；静态 I/Q--STFT 融合接近 CLDNN 但有整体损失；标量门控融合低于静态融合，构成负结果；MCLDNN 两个训练种子单类塌陷被保留登记。提案模型 \path{fusion_cldnn_stft} 把 fusion 的 I/Q 分支升级为 CLDNN 风格编码器，叠加保标签的相位旋转加循环时移信号域增强（仅训练）；在同一固定划分与三训练种子下取得平均整体准确率 0.6284（\(\sigma=0.0011\)），相对 CLDNN 基线提升 1.55 个百分点，且整体、低、中、高信噪比四项指标同时正向；avg(\(\mathrm{SNR}>0\)) 维度数值上高于 MCLDNN 原论文报告的 0.9112。第~4.7.4~节的单干预消融进一步显示信号域增强是主要涨点来源；额外叠加标签平滑 0.1 在本数据集上按论文自身的配对统计标准显著有害（McNemar $p = 0.0155$），因此最终提案模型不包含标签平滑。提案模型与 2024 年代复值多流／Transformer 混合基线量级一致，与 LENet-M（0.6463）和 SigFormer（0.6371）仍有 1--2 个百分点差距。本研究受单数据集、硬件变更与未做单干预消融三条边界限制；跨数据集一致性留作未来工作。
\end{cnabstract}
